\begin{figure}[h]
	\centering
	\begingroup

\definecolor{mscXEvent}{HTML}{D62728}
\definecolor{mscYEvent}{HTML}{009E73}
\definecolor{mscAfterT}{HTML}{E69F00}

\tikzset{
	msc x wire/.style={msc wire above=mscXEvent},
	msc y wire/.style={msc wire above=mscYEvent},
	msc after t wire/.style={msc wire above=mscAfterT},
	msc x control/.style={msc control left=mscXEvent},
	msc circle event/.style={
		msc event,
		shape=yquant-circle,
	},
	msc x event/.style={
		msc circle event,
		fill=mscXEvent!25,
		draw=mscXEvent!85!black,
	},
	msc y event/.style={
		msc circle event,
		fill=mscYEvent!25,
		draw=mscYEvent!85!black,
	},
}

\yquantset{
	msc x cnot/.style={
		every control line/.append style={msc x control},
	},
}

\begin{tikzpicture}
\begin{yquant}

qubit {0} q0;
qubit {a} qa;
qubit {b} qb;
qubit {1} q1;
qubit {c} qc;
qubit {2} q2;
qubit {d} qd;
qubit {3} q3;

[msc x event]
box {X} q1;
setstyle {msc x wire} q1;
align q0,qa,qb,q1,qc,q2,qd,q3;

[msc x cnot]
cnot qa|q1;
setstyle {msc x wire} qa;
cnot qd|q2;

[msc x cnot]
cnot q0|qa;
setstyle {msc x wire} q0;
cnot q3|qd;
cnot q2|qc;
cnot q1|qb;

[msc y event]
box {Y} q3;
setstyle {msc y wire} q3;
align q0,qa,qb,q1,qc,q2,qd,q3;

dmeter {X} qa,qb,qc,qd;
discard qa,qb,qc,qd;
box {T} q0;
box {T$^\dag$} q1;
box {T} q3;
box {T} q2;

setstyle {msc after t wire} q0,q1,q3;
hspace {1mm} q0,qa,qb,q1,qc,q2,qd,q3;

\end{yquant}
\end{tikzpicture}
\endgroup

	\hfill
	\includegraphics[width=0.77\linewidth]{figures/summary.pdf}
	\begin{subfigure}[b]{0.24\linewidth}
		\centering
		\caption{In T-state cultivation,
		two Pauli error events can propagate into a weight-3 Clifford error
		that can lead to an undetected logical error.}
		\label{fig:2_fault_configuration}
	\end{subfigure}
	\hfill
	\begin{subfigure}[b]{0.27\linewidth}
		\centering
		\caption{The double check fault distance
		is 2 (3) for distance-3 (-5) T-state cultivation,
		rather than 3 (5) for its S-state proxy.
		Missing datapoints are exactly zero.}
		\label{fig:lep_against_fault_count}
	\end{subfigure}
	\hfill
	\begin{subfigure}[b]{0.43\linewidth}
		\centering
		\caption{In distance-5 T-state cultivation,
		detecting the degrading hook errors with our Z flags
		improves the simulated pre-escape logical error rate per kept shot.
		Shading shows the region within a factor of 1000 of the maximum likelihood.}
		\label{fig:d5_pre_escape_ler}
	\end{subfigure}
	\caption{Summary of results.
	Blue (orange) plot data correspond to cultivating S (T) states.}
	\label{fig:summary}
\end{figure}

\section{Introduction}
\label{sec:introduction}

Magic state cultivation (MSC)~\cite{Gidney2024},
or T-state cultivation,
is a relatively new protocol for producing logical T states
that has recently exploded in popularity
\cite{Chen2026,Claes2025,Sahay2026,Vaknin2025,Wan2025b,Hirano2025a,Rosenfeld2025,Hetenyi2026,Chen2026a}
due to its practicality:
it shrinks the spacetime cost of a logical T state
to that of a lattice surgery CX.
Simulating MSC on a classical computer is costly
due to its non-Clifford nature,
so early numerical benchmarks \cite{Xu2025,Ibe2025,Hines2026}
relied on simulating a Clifford proxy,
S-state cultivation,
on a stabiliser simulator such as Stim~\cite{Gidney2021a}.
Since the T- and S-state cultivation circuits share the same structure,
it was expected that their logical error rates would be similar.
However, simulations from the original paper~\cite{Gidney2024} showed that
distance-3 T-state cultivation in its early stages
has a (roughly 2$\times$) higher logical error rate
than its S-state analogue.
Subsequent papers \cite{Wan2026,Chen2026b} observed that
this T--S discrepancy widens as the noise level decreases.
More recent papers \cite{Li2025,Tuloup2026,Chase2026,Fang2026a}
simulate distance-5 cultivation
and observe an even larger (roughly 11$\times$) T--S discrepancy
that also widens as the noise level decreases;
this has yet to be explained.
This raises concerns for resource estimates that rely on MSC,
and highlights the danger of
using Clifford proxies to simulate
the error behaviour of non-Clifford protocols.

Motivated by this,
we analytically explain the widening T--S discrepancy
with hopes to improve cultivation in the future.
We adopt the approach taken by Daguerre and Kim~\cite[\S V]{Daguerre2025},
which is to consider the Clifford errors that result from
propagating Pauli errors through a layer of non-Clifford gates.
Then,
instead of Monte Carlo sampling,
we exhaustively enumerate the dominant contributions
to the logical error rate,
which directly tells us the \emph{fault distance}
i.e.\ level of logical-error suppression in the physical noise level.
It is known that the fault distance of a subcircuit,
called the \emph{double check},
within distance-3 (-5) S-state cultivation is 3 (5): as expected.
To our knowledge,
we are the first to show that it is instead 2 (3) in T-state cultivation;
see \cref{fig:lep_against_fault_count}.
This is due to the following reason:
the CXs spread Pauli errors which are then rotated
to multiqubit Clifford errors by the physical T-gates in the circuit;
\cref{fig:2_fault_configuration} shows an example.
Rather than merely changing the signs of the code stabilisers,
such an error transforms the code into a different stabiliser code
whose codespace can overlap the original,
allowing postselection to accept a logically erroneous component.
By contrast,
the corresponding errors remain Pauli in S-state cultivation
and are deterministically detected.

To analyse these Clifford errors,
we derive a formula (\cref{theorem:probability_trivial_syndrome})
for the probability to observe a trivial syndrome
for any mixed state encoded in any Pauli stabiliser code afflicted with any Clifford error.
We also provide an algorithm (\cref{alg:probability_trivial_syndrome})
to compute this probability in $\mathcal O(n^3)$ time
for fixed logical-qubit count,
where $n$ is the physical qubit count.
These two general results may be of independent interest.
We then calculate the fidelity (\cref{prop:resulting_fidelity})
specifically for T states afflicted with
the Clifford errors that are encountered in MSC,
conditioned on a trivial syndrome.

Having enumerated the malignant Clifford errors in MSC,
we use the technique of flags~\cite{Chao2018}
to design 2D-local Z-flagged versions of the double checks that detect those errors.
We then use the fast universal simulator SymFT~\cite{Fang2026a}
to simulate the whole MSC protocol except
the final so-called `escape' stage, which is purely Clifford.
Monte Carlo sampling shows these Z flags recover
the intended scaling of logical error rate with noise level,
which suggests hook Clifford errors
are the sole reason for the degraded fault distance of pre-escape MSC.
These sampling results are summarised in \cref{fig:d5_pre_escape_ler}
which shows the practical improvement
these Z flags bring to distance-5 T-state cultivation.
This demonstrates that the fault tolerance of non-Clifford circuits
can be improved by considering the propagation of non-Pauli errors,
even with the simple addition of more Clifford gates.

\Cref{sec:background} covers the prerequisite theory.
We first show the lower-than-expected fault distance in MSC
in \cref{sec:fault_distance_of_magic_state_cultivation},
then present our general Clifford-error analysis in \cref{sec:clifford_error_analysis}.
We quantitatively analyse the fault configurations in \cref{sec:thorough_analysis}
before evaluating our Z-flagged version of MSC in \cref{sec:z_flagged_double_checks}
and concluding in \cref{sec:conclusion}.
\Cref{sec:formal_definitions} provides definitions
required for the rest of the appendices
and \cref{sec:additional_figures} shows alternative figures to those in the main text.
The code for \cref{alg:probability_trivial_syndrome}
and for all the numerics in this paper
is on GitHub at \cite{Chan2026c_quantum_bibstyle}.